%Explicaci\'on del problema que se va a estudiar, en qu\'e consiste, cu\'ales son sus variables , restricciones y objetivo(s) de manera general (en palabras, no una formulaci\'on matem\'atica). Debe entenderse claramente el problema y qu\'e busca resolver.
%Explicar si existen problemas relacionados.
%Destacar, si existen, las variantes m\'as conocidas.\\
%Redactar en tercera persona, sin faltas de ortograf\'ia y referenciar correctamente sus fuentes mediante el comando  \verb+\cite{ }+. Por ejemplo, para hacer referencia al art\'iculo de algoritmos h\'ibridos para problemas de satisfacci\'on
%de restricciones~\cite{Prosser93Hybrid}.

\subsection{Contexto}

El \textit{Truck Drone Parcel Delivery Problem} (TDPDP) es un problema donde se busca encontrar una ruta para un camión de entregas, el cual debe pasar por varios puntos de entrega. El TDPDP tiene la adición de un dron, el cual permite hacer entregas con tiempos distintos a los del camión.

\subsection{Definición Formal}

Para poder entender el problema del TDPDP, primero se muestra que es una variante de otros problemas más conocidos en el mundo de la optimización. Estos corresponden al TSP (\textit{Traveling Salesperson Problem}, o el Problema del Vendedor Viajero en español) y al problema de VRP (\textit{Vehicle Routing Problem}, o el Problema de Ruteo de Vehículos en español). El VRP es la variante más cercana al TDPDP, y la cual es también variante del TSP, por lo que se explica el VRP para poder extenderlo al TDPDP más adelante.

En el VRP, definido según \textit{Irnich et al.} \cite{Irnich2014FamilyVRP}, se tiene que \textbf{dado} un conjunto de entregas solicitadas (también entendidas como ubicaciones) y una cantidad de vehículos, la \textbf{tarea} es determinar una ruta para cada uno de los vehículos para poder llegar a las ubicaciones de entrega a un costo mínimo, es decir, decidir qué vehículo va a cuáles ubicaciones y en qué secuencia para que todas (o algunas) de las ubicaciones sean cubiertas. 

De la definición del VRP, se define una variante conocida como CVRP (\textit{Capacitated Vehicle Routing Problem}), la cual se acerca más a la definición adoptada para el TDPDP. Esta variante tiene un punto de inicio (conocido como \textit{depot}, depósito en español), y se hace también referencia a puntos (o ubicaciones) de entrega. El problema debe resolverse desde el punto de inicio definido hacia cada una de las ubicaciones de entrega. Las restricciones del problema siguen al VRP, es decir, se tiene la restricción de visitar todas las ubicaciones de entrega, pero se le incluye la restricción de empezar desde el depósito y volver al depósito.

Conociendo el problema de VRP, el TDPDP es una variante del VRP en la que se incluye un dron, el cual permite realizar entregas en conjunto con el vehículo. Siguiendo la definición realizada por \textit{Murray} y \textit{Chu} \cite{Murray2015Flying}, este dron tiene las características de que su costo para viajar desde un nodo puede ser distinto al costo del camión y, adicionalmente, puede hacer entregas en paralelo mientras el camión visita otros nodos, lo cual permite hacer el costo asociado al viaje más bajo y los tiempos de entrega más cortos. 

\subsection{Restricciones}

Adicionalmente a las restricciones del VRP, se tienen restricciones exclusivas del TDPDP relacionadas con el dron. Según \textit{Murray} y \textit{Chu} \cite{Murray2015Flying}, hay restricciones con respecto a la cantidad de drones por camión, la batería del dron y las velocidades de vuelo. Solo puede haber un dron asociado a cada camión, el cual al salir del camión debe reunirse con este en algún nodo posterior. El dron tiene una batería asociada, la cual indica cuánta distancia puede recorrer y cuánto tiempo puede quedarse en espera (por ejemplo, un dron podría esperar en un nodo al camión siempre y cuando le conserve batería suficiente cuando este llegue), y por último la velocidad del dron limita qué nodos pueden visitarse.

\subsection{Objetivo}

El objetivo del TDPDP es minimizar el tiempo de retorno más tardío de ambos vehículos al depósito, lo cual se traduce en determinar cuál es la ruta que se realiza. 

Las soluciones del problema vienen dadas según la selección de qué nodos se visitan con el camión, cuáles se visitan con el dron y en qué orden se realizan estas entregas. La ventaja de este modelo en comparación con un VRP es el tiempo que se puede ahorrar utilizando el dron para entregas más rápidas que las realizadas con el camión.

Se proponen otras variantes para la función objetivo. En \textit{Thomas et al.} \cite{thomas2024collaborative} se menciona minimizar el costo operacional, donde el costo se asocia a la distancia recorrida por dron y camión, además del tiempo ocioso del camión. También en \textit{Moshref-Javadi} y \textit{Winkenbach} \cite{moshrefjavadi2021applications} se propone medir como función objetivo la suma de tiempos de espera para cada cliente con su pedido.

\subsection{Complejidad}

Se tiene evidencia respecto a que el problema de TDPDP es \textit{NP-hard}. El problema de VRP es definido como \textit{NP-hard} en \textit{Irnich et al.} \cite{Irnich2014FamilyVRP}, y al ser el TDPDP una variante del VRP, se entiende que el problema es igualmente \textit{NP-hard}. Este fundamento es desarrollado por \textit{Ponza} \cite{ponza2016optimization} mediante una reducción al problema del TSP, al cual el VRP también se reduce. El análisis de complejidad se hace más profundo en \textit{Boysen et al.} \cite{boysen2018drone}, en donde se fundamenta que el problema se mantiene \textit{NP-hard} para instancias en donde se varía el número de drones o donde se tiene al camión con una ruta prefijada.

Dado que el problema es \textit{NP-hard}, se concluye que no existen algoritmos exactos de tiempo polinomial conocidos para encontrar una solución óptima en instancias generales.

\subsection{Variantes}

Se presenta una tabla con las principales variantes del problema, variando generalmente en la forma en que el camión y el dron distribuyen la carga. Se profundizará en algunas de estas interpretaciones en el análisis del estado del arte.

\begin{table}[H]
\centering
\scriptsize
\setlength{\tabcolsep}{3.5pt}
\caption{Principales variantes del modelo TDPDP}
\label{tab:variantes-tdpdp}
\begin{tabular}{p{2.3cm}p{3.4cm}p{7.9cm}}
\hline
\textbf{Modelo} & \textbf{Referencia} & \textbf{Características} \\
\hline
FSTSP & \textit{Murray} y \textit{Chu} (2015) \cite{Murray2015Flying} & Un camión y un dron colaboran; el dron es lanzado desde un nodo visitado por el camión, realiza una entrega y es recuperado posteriormente en otro punto de la ruta. \\
\hline
PDSTSP & \textit{Murray} y \textit{Chu} (2015) \cite{Murray2015Flying} & Uno o varios drones operan desde el depósito en paralelo al camión, realizando entregas directas sin necesidad de sincronización directa con la ruta terrestre. \\
\hline
FSTSP (OPL / CPLEX) & \textit{Ponza} (2016) \cite{ponza2016optimization} & Reformulación con variables de tiempo que indican cuándo está listo para partir ($t^T_i, t^D_i$), optimizando la cota de autonomía de batería del dron. \\
\hline
TSP-D & \textit{Agatz et al.} (2018) / \textit{Ponza} (2016) \cite{ponza2016optimization} & Variante flexible que permite configuraciones donde los nodos de lanzamiento y recuperación pueden coincidir. \\
\hline
DSP-A & \textit{Boysen et al.} (2018) \cite{boysen2018drone} & Programación de drones con ruta de camión fija bajo Política A (despegue y retorno en la misma parada). Polinomial para 1 dron y \textit{NP-hard} para múltiples drones. \\
\hline
DSP-B & \textit{Boysen et al.} (2018) \cite{boysen2018drone} & Programación con ruta fija bajo Política B (retorno a lo más en la parada sucesiva). Resoluble en tiempo polinomial mediante reducción al \textit{Linear Assignment Problem} (LAP). \\
\hline
DSP-C & \textit{Boysen et al.} (2018) \cite{boysen2018drone} & Programación con ruta fija bajo Política C (retorno irrestricto en cualquier parada futura). \textit{NP-hard} por reducción. \\
\hline
CTDRSP-ELR & \textit{Thomas et al.} (2024) \cite{thomas2024collaborative} & Permite que el dron pueda volver al camión si este está en marcha. Incluye despliegue en parada fija, despliegue en movimiento, recuperación en movimiento y recuperación en parada fija. \\
\hline
Modelo PD & \textit{Moshref-Javadi} y \textit{Winkenbach} (2021) \cite{moshrefjavadi2021applications} & Reparto con drones desde depósitos fijos hacia clientes sin necesidad de un camión. Limitado por la capacidad de batería del UAV y la congestión en el almacén. \\
\hline
\end{tabular}
\end{table}
